%% ============================================================
%% pulsar/unforgeability.tex
%% ============================================================
%%
%% Canonical Pulsar SUF-CMA unforgeability theorem.
%% Papers cite this; do not redefine.
%% ============================================================

\subsection*{Pulsar threshold signature unforgeability}

\begin{theorem}[Pulsar SUF-CMA under MLWE/MSIS]
\label{thm:pulsar-suf-cma}
Let $\mathcal{A}$ be a probabilistic polynomial-time adversary against
Pulsar with at most $q_S$ signing queries and at most $t-1$ corrupted
parties. Then
\[
  \mathrm{Adv}^{\mathsf{SUF\text{-}CMA}}_{\mathsf{Pulsar}}(\mathcal{A})
  \;\leq\;
  q_S \cdot \mathrm{Adv}^{\mathsf{MLWE}}_{q, \sigma_E}(\mathcal{B}_1)
  \;+\;
  \mathrm{Adv}^{\mathsf{MSIS}}_{q, [\mathbf{A}|-c\tilde{\mathbf{b}}], 2B}(\mathcal{B}_2)
  \;+\;
  \mathrm{negl}(\lambda).
\]
The reduction algorithms $\mathcal{B}_1$ (MLWE distinguisher) and
$\mathcal{B}_2$ (MSIS solver) run in time polynomial in $\mathcal{A}$'s
running time. Concrete parameters in
\texttt{quasar-cert-soundness.tex} App.~E give classical security
$\geq 2^{142}$ and quantum security $\geq 2^{130}$ via BDGL sieving
plus Grover speedup.

\textbf{Lean mechanization:} \texttt{Crypto.Ringtail.\\
ringtail\_pq\_unforgeability} in \texttt{lean/Crypto/Ringtail.lean};
\texttt{Crypto.Threshold.Reshare.reshare\_preserves\_public\_key} in
\texttt{lean/Crypto/Threshold/Reshare.lean}.
\end{theorem}

\begin{remark}[Hash-suite invariance]
\label{rem:hash-suite-invariance}
The reduction is independent of the choice of
$\mathsf{HashSuiteID} \in \{\mathsf{Pulsar\text{-}SHA3},
\mathsf{Pulsar\text{-}BLAKE3}\}$ provided the suite is collision-resistant
and behaves as a programmable random oracle in the SUF-CMA game. Both
shipped suites satisfy this; activation certificates and signatures
under different suites are distinct objects (the suite identifier is
bound into the transcript, Definition~\ref{def:pulsar-transcript}).
\end{remark}
